% ██████████████████████████████████████████████████████████████████
%  CHAPTER 4 — SYSTEM ANALYSIS AND DESIGN
% ██████████████████████████████████████████████████████████████████
\chapter{System Analysis and Design}
\label{ch:system_design}

\section{Functional Requirements}
\label{sec:design_functional}

\begin{enumerate}[label=\textbf{FR\arabic*:}]
  \item The system shall provide late-payment risk prediction for any open invoice owned by an authenticated tenant.
  \item The system shall provide a SHAP-based feature attribution for each prediction served via the local-prediction tier.
  \item The system shall determine the appropriate service tier (local-prediction vs.\ global-heuristic) automatically, based on the tenant's historical invoice volume.
  \item The system shall expose the tier decision and its supporting count to API consumers as first-class fields of the prediction response.
  \item The system shall persist global heuristics in a versioned JSON artefact loadable at request time.
\end{enumerate}

\section{Non-Functional Requirements}
\label{sec:design_nonfunctional}

\begin{enumerate}[label=\textbf{NFR\arabic*:}]
  \item \emph{Backward compatibility.} The introduction of the tier gate shall not alter the signature or return type of the existing \texttt{predict()} method.
  \item \emph{Auditability.} The tier decision shall be deterministic given (tenant, threshold, artefact version) and shall be logged via the existing prediction-logging infrastructure.
  \item \emph{Performance.} The heuristic JSON shall be memoised on first access, with a worst-case file-system read incurred at most once per process lifetime.
  \item \emph{Tenant isolation.} The tier gate shall not leak cross-tenant invoice data; the invoice count is scoped strictly to the requesting company.
\end{enumerate}

\section{Conceptual Architecture}
\label{sec:design_concept}

The Maturity-Aware AI Routing pattern introduces a thin gating layer between the prediction-request entry point and the existing trained-model artefact. The gating layer is responsible for three concerns: (i) computing the tenant's data-maturity signal, (ii) selecting the appropriate service tier, and (iii) loading the cold-start heuristic when required.

\begin{figure}[H]
  \centering
  \fbox{\parbox{0.85\textwidth}{\centering\vspace{1.5cm}
    \textit{[Placeholder for Figure~\ref{fig:routing_pattern}]}\\[0.4em]
    Component diagram: \emph{Prediction Request} $\rightarrow$
    \emph{Tier Gate} $\rightarrow$ \{\emph{Local Predictor} $|$
    \emph{Global Heuristic Loader}\} $\rightarrow$
    \emph{Unified Response Object}.
    \vspace{1.5cm}}}
  \caption{Conceptual architecture of the Maturity-Aware AI Routing pattern.}
  \label{fig:routing_pattern}
\end{figure}

\section{Data Model}
\label{sec:design_data}

The tier gate consumes the existing \texttt{Invoice} model without schema changes; the maturity signal is derived from a filtered count of \texttt{status~$\in$~\{paid, overdue, partial\}}. The global heuristic is persisted as a flat JSON document containing (i) a human-readable \texttt{heuristic} narrative string and (ii) a ranked list of \texttt{top\_features} with normalised importance weights.

\section{Sequence of Operations}
\label{sec:design_sequence}

\begin{enumerate}[label=\arabic*.]
  \item An API consumer requests a prediction for a specific invoice.
  \item The view-layer instantiates \texttt{LatePaymentPredictor} bound to the requesting company.
  \item The predictor's \texttt{predict\_or\_heuristic} method computes the tenant's qualifying invoice count.
  \item If the count is below \texttt{COLD\_START\_INVOICE\_THRESHOLD}, the heuristic JSON is loaded (or fetched from memoised cache) and returned with \texttt{mode = "global\_heuristic"}.
  \item Otherwise, the local model is invoked via \texttt{predict()} and the response is returned with \texttt{mode = "local\_prediction"}.
  \item The response is serialised to JSON and returned to the consumer.
\end{enumerate}
